\documentclass[11pt]{article}
\usepackage{amsmath}
\usepackage{amssymb}
\usepackage{amsthm}
\usepackage{enumitem}
\usepackage[margin=1in]{geometry}
\usepackage{multicol}
\usepackage{tikz}

% Load hyperref to automatically generate the PDF outline from structural tags
\usepackage[pdftex, bookmarks=true, bookmarksnumbered=true]{hyperref}

% Optional: Use the bookmark package for faster rendering and advanced control
\usepackage{bookmark} 

\newenvironment{case}[2]{\par\textbf{Case #1: #2.}}{}
\title{Homework \#2\\{\large Math 104: Intro to Analysis \\ Summer 2026} \\ {\normalsize UC Berkeley}}
\date{July 10, 2026}
\author{Donald Aingworth IV}

\begin{document}
\renewcommand\qedsymbol{OE$\Delta$}
\newcommand{\abs}[1]{\left| #1 \right|}
\newcommand{\andtxt}{\text{ and }}
\newcommand{\Ais}{\ensuremath{\overset{A}{=}}}
\newcommand{\Cis}{\ensuremath{\overset{C}{=}}}
\newcommand{\closure}[1]{\ensuremath{\overline{#1}}}
\newcommand{\encircle}[1]{%
  \tikz[baseline=(char.base)]{%
    \node[shape=circle, draw, inner sep=1pt] (char) {#1};%
  }%
}
\newcommand{\mathif}{\text{if }}
\newcommand{\seq}[2][n]{\ensuremath{\left\{ #2 _{#1} \right\}}}
\newcommand{\zz}[1]{$\mathbb{Z}/ #1 \mathbb{Z}$}

\maketitle

\tableofcontents

\pagebreak
\setcounter{section}{-1}
\section{Various Theorems}
    Various theorems, propositions, and statements that will be useful in this problem set will be detailed here.
    \subsection{Reverse Triangle Inequality}
        In a metric space $(X,d)$ with elements $a, b, c \in X$, $\abs{d(a,b) - d(b,c)} \leq d(a,c)$.
        \begin{proof}
            Let $a, b, c \in X$.
            By the triangle equality, we know $d(a,b) \leq d(a,c) + d(b,c)$ and $d(b,c) \leq d(a,c) + d(a,b)$.
            Subtracting $d(a,b)$ from one and $d(b,c)$ from the other, we find $d(a,b) - d(b,c) \leq d(a,c)$ and $d(b,c) - d(a,b) \leq d(a,c)$.
            Reversing the latter, we find $d(a,b) - d(b,c) \geq -d(a,c)$.
            This means that $-d(a,c) \leq d(a,b) - d(b,c) \leq d(a,c)$.
            Ergo, $\abs{d(a,b) - d(b,c)} \leq d(a,c)$.
        \end{proof} 

    \subsection{Specific for all means general exists}
        Suppose $A \subset B$ and $A \neq \varnothing$ and $X$ refers to some statement that can be either true or false. If $X(a)$ is true for all $a \in A$, then there exists some $b \in B$ where $X(b)$ is true.
        \begin{proof}
            Let $X(a)$ be true for all $a \in A$.
            Take any $a \in A$.
            Since $A \subset B$, $a \in B$.
            Therefore, there exists some $a \in B$ where $X(a)$ is true.
        \end{proof}
        Used in 3a.

    \subsection{Monotone decreasing sequences converge to Infimum}
        If a sequence is monononically decreasing, then it converges to its infimum.
        \begin{proof}
            Suppose a monotonically decreasing sequence $s_n$, so for all $n \in \mathbb{N}$, $s_n \geq s_{n + 1}$.
            Suppose further that it is bounded below by $s = \inf\{s_n\}$, so for any $n \in \mathbb{N}$, $s_n > s$.

            Let $\varepsilon > 0$.
            There is $N \in \mathbb{N}$ such that $s_N < s + \varepsilon$.
            Then, for all $n > N$, $s + \varepsilon > s_N \geq s_n > s$.
            Therefore, $\varepsilon > s_n - s > 0$.
            Therefore, $\varepsilon > \abs{s_n - s}$.
            Ergo, $\{s_n\}$ converges to $\inf\{s_n\}$.
        \end{proof}

    \subsection{Minimum is Infimum}
        If for all $x \in X$, $y \leq x$ and $y \in X$, then $y = \inf X$.
        \begin{proof}
            Since $y \leq x$ for all $x \in X$, $y$ is a lower bound of $X$.
            Let $\varepsilon > 0$.
            Therefore, $y + \varepsilon > y$.
            Since $y \in X$, for all $\varepsilon > 0$, there exists some element in $x \in X$ such that $x < y + \varepsilon$.
            Ergo, $y = \inf X$.
        \end{proof}
        A similar proof can be done to prove that if for all $x \in X$, $y \geq x$ and $y \in X$, then $y = \sup X$.\\
        Used in 3c.
    
    \subsection{Floor and Ceiling Functions}
        \textbf{Floor function ($\left\lfloor \right\rfloor$):}
            $n = \left\lfloor x \right\rfloor$ is defined as the integer $n$ such that $n \leq x < n + 1$.\\
        \textbf{Ceiling function ($\left\lceil \right\rceil$):}
            $n = \left\rceil x \right\rceil$ is defined as the integer $n$ such that $n - 1 < x \leq n$.

    \subsection{$2^n > n$}
        If $n \in \mathbb{N}$, then $2^n > n$.
        \begin{proof}[Proof by Induction]
            \textbf{Base case:}
                We know that $2^1 = 2 > 1$.

            \textbf{Inductive step}
                Suppose that $2^n > n$.
                Furthermore, since $n \geq 1$, $2n \geq n + 1$.
                Multiply both sides of the inequality by $2$ to find $2^{n+1} > 2n$.
                Transistively, this means that $2^{n+1} > n + 1$.

            Ergo, if $n \in \mathbb{N}$, then $2^n > n$.
        \end{proof}
        Used in 2c.

\pagebreak
\section{Problem 1}
    Let $(X, d)$ be a metric space. A subset $E \subset X$ is called \textit{dense} in $E$ if, for any open $O \subset X$, we have $O \cap E \neq \varnothing$. A metric space $(X, d)$ is called \textit{separable} if it has a countable dense subset.
    \begin{enumerate}[label=(\alph*)]
        \item 
        Show that $\mathbb{Q}$ is dense in $\mathbb{R}$, and conclude that $\mathbb{R}$ is separable.

        \item
        Show that a subset $E$ is dense in $X$ if and only if $\overline{E} = X$.

        \item 
        By following these steps, prove that any compact metric space $(X, d)$ is separable:
        \begin{enumerate}[label=(\roman*)]
            \item 
            Using problem 5(c) on homework 2, show that for each $n \in \mathbb{N}$, the following is true: there is a finite set $S_n \subset X$ so that every point of $X$ is within distance $\frac{1}{n}$ of some point of $S_n$.

            \item 
            Let $S = \bigcup\limits_{n = 1}^{\infty} S_n$. Show that $S$ is countable.

            \item 
            Show that every point of $X$ is either in $S$ or a limit point of $S$. Conclude that $S$ is dense in $X$, and hence that $X$ is separable.
            
        \end{enumerate}
        
    \end{enumerate}
    % \pagebreak
    \subsection{Solution (a)}
        \textbf{Alexis's Proposition:}
            For any $a, b \in \mathbb{R}$ where $a < b$, there exists some $\frac{m}{n} \in \mathbb{Q}$ such that $a < \frac{m}{n} < b$.
            \begin{proof}
                Let $a, b \in \mathbb{R}$ where $a < b$.
                Since $a < b$, $b - a > 0$.
                By the Archimedean property, there exists some $n \in \mathbb{N}$ such that $n(b-a)>1$.
                Distribute and add $na$ to both sides to find $nb > na + 1$.
                Also by the Archimedean property, there exists some natural numbers $m_1$ and $m_2$ such that $m_1 > na$ and $m_2 > -na$.
                The latter means that $na > -m_2$.
                Therefore, $m_1 > na > -m_2$.
                Since $m_1$ and $-m_2$ are both integers, there exists some integer $m$ such that $-m_2 < m < m_1$ and $m - 1 \leq na < m$.
                This is equivalent to $m \leq na + 1 < m + 1$.
                Combine all of these and we find that $na < m \leq na + 1 < nb$.
                Shortening this, $na < m < nb$.
                Divide this all by $n$ to find $a < \frac{m}{n} < b$.
                Since $n$ is a natural number and $m$ is an integer, $\frac{m}{n}$ is a rational number.
                Ergo, if $a, b \in \mathbb{R}$ and $a < b$, there exists some $\frac{m}{n} \in \mathbb{Q}$ such that $a < \frac{m}{n} < b$.
            \end{proof}

            Now we prove $Q$ is dense in $\mathbb{R}$ with the usual metric.
            \begin{proof}
                Take any set $O$ open in $\mathbb{R}$.
                This means that for any $o \in O$, there exists some $r > 0$ such that there exists an open ball $B_r(o)$ that is a subset of $O$.
                We define this open ball as $B_r(o) = \left\{ x \in \mathbb{R} \Big| \abs{x - o} < r \right\}$.

                Since $r > 0$, this means $r > -r$ and therefore $o + r > o - r$.
                By Alexis's Proposition, there exists some rational number $\frac{m}{n}$ such that $o - r < \frac{m}{n} < o + r$.
                Subtract $o$ from all sides to find $-r < \frac{m}{n} - o < r$.
                By the definition of the absolute value, $\abs{\frac{m}{n} - o} < r$.
                Therefore, $\frac{m}{n} \in B_r(o)$ and by extension $\frac{m}{n} \in O$.
                Therefore, $\frac{m}{n} \in \mathbb{Q} \cap O$.
                Therefore, $\mathbb{Q} \cap O \neq \empty$.
                This means that for any open $O \subset \mathbb{R}$, $\mathbb{Q} \cap O \neq \empty$.
                Ergo, $\mathbb{Q}$ is dense in $\mathbb{R}$.
            \end{proof}

            Last we prove that $\mathbb{R}$ is seperable.
            \begin{proof}
                We know that $\mathbb{Q}$ is dense in $\mathbb{R}$.
                We know that $\mathbb{Q}$ is a subset of $\mathbb{R}$.
                We know that $\mathbb{Q}$ is countable.
                Therefore, $\mathbb{Q}$ is a countable, dense subset of $\mathbb{R}$.
                Ergo, $\mathbb{R}$ is seperable.
            \end{proof}

    \pagebreak
    \subsection{Solution (b)}
        First, we prove that if $E$ is dense in $X$, then $\closure{E} = X$.
        \begin{proof}
            Suppose that $E$ is dense in $X$.
            Take any point $x \in X$.
            Create an open ball of radius $r$ for any $r > 0$.
            We can define the open ball as $B_r(x) = \left\{ y \in X \Big| \abs{y - x} < r \right\}$.
            Since open balls are open and $E$ is dense in $X$, there is a point $p \in E \cap B_r(x)$.
            Therefore, there is a point $p$ in $E$ in every open ball centered at any point in $X$.
            Therefore, every point in $X$ is a limit point of $E$.
            Therefore, $X$ is the set of limit points of $E$
            Since \closure{E} is the union of $E$ and the set of its limit points and every point in $E$ is in $X$ (since $E$ is a subset of $X$), $X = \closure{E}$.
        \end{proof}

        Next we prove that $E$ is dense in $X$ if $\closure{E} = X$.
        \begin{proof}
            Let $\closure{E} = X$.
            Take any open set $O \subset X$.
            Since $\closure{E} = E \cup E'$, for any point $p \in X$, either $p \in E$ or $p \in E'$.

            If $p \in E$, then $p \in E \cap O$, so $E \cap O \neq \varnothing$.

            If $p \in E'$, then $p$ is a limit point of $E$.
            Since $O$ is open, there is some open ball of radius $r$ centered around $p$ such that $B_r(p) \subset O$.
            Since $p$ is a limit point of $E$, there exists some $q \in E$ such that $p \neq q$ and $q \in B_r(p) \subset O$.
            Therefore, since $q \in E$ and $q \in O$, $q \in O \cap E$.
            Therefore, $E \cap O \neq \varnothing$.

            Ergo, if $\closure{E} = X$, $E$ is dense in $X$.
        \end{proof}

    \pagebreak
    \subsection{Solution (c)}
        \subsubsection{Part (i)}
            \begin{proof}
                We know that $X$ is compact.
                From Problem 5c on Homework 2, we know that for a compact set has the following property: for all $\varepsilon > 0$, there is a finite set $\{x_1, \dots, x_k\} \subset X$ such that each $y \in X$ is within distance $\varepsilon$ of some $x_j$.
                Suppose that $\varepsilon = \frac{1}{n}$ for some $n \in \mathbb{N}$.
                Label the finite set for this value of $n$ $\{x_1, \dots, x_k\}$ as $S_n$.
                Ergo, for any $n \in E$, there is a finite set $S_n \subset X$ so that every point of $X$ is within distance $\frac{1}{n}$ of some point in $S_n$. 
            \end{proof}
        \subsubsection{Part (ii)}
            \begin{proof}
                We know that for any $n \in \mathbb{N}$, $S_n$ is finite.
                We also know that the union of at most countable sets of at most countable cardinality is at most countable.
                Therefore, since we are working with countable sets that are finite, the union of $S_n$ for all $n \in \mathbb{N}$ is countable.
                Ergo, $S = \bigcup\limits_{n = 1}^{\infty} S_n$ is countable.
            \end{proof}
        \subsubsection{Part (iii)}
            \begin{proof}
                We have two cases here: for any $x \in X$, either $x \in S$ or $x \notin S$.
                If $x \in S$, then we don't need anything further.

                Suppose instead that $x \notin S$.
                We proved earlier that every point $x \in X$ is within distance $\frac{1}{n}$ of some point $s \in S$ for any $n \in \mathbb{N}$, or in other words $d(x,s) < \frac{1}{n}$.
                For $x$ to be a limit point of $S$, it must be that for any $r > 0$, there is a point $y \in S$ such that $y \neq x$ and $y \in B_r(x)$ for the open ball of radius $r$ centered at $x$.
                $B_r(x)$ is defined as $B_r(x) = \left\{ z \in X \Big| d(x,z) < r \right\}$.
                The Archimedean property proved that for any $a > 0$, there exists some natural number $n$ such that $a \geq \frac{1}{n} > 0$.
                Therefore, for any $r > 0$, there exists some natural number $n$ such that $r \geq \frac{1}{n} > 0$.
                Furthermore, we know from earlier that for ever point $x \in X$, there is some point $s \in S$ such that $d(x,s) < \frac{1}{n}$ for any $n \in \mathbb{N}$.
                Transistively, this means that for any $r > 0$ and for any $x \in X$, there is some $s \in S$ such that $d(x,s) < r$.
                Since $x \notin S$, $x \neq s$, so $d(x,s) > 0$.
                Therefore, $s \in B_r(x)$ and $s \in S$ for any $r > 0$ and for any $x \in X \backslash S$.
                Therefore, $x$ is a limit point of $S$.
                
                Ergo, every point of $X$ is either in $S$ or is a limit point of $S$.
                This means that $X = \closure{S}$.
                This means that $S$ is dense in $X$ by part (ii).
                Sinc $S$ is countable and a metric space is countable if it has a dense subset, $(X,d)$ is seperable.
            \end{proof}

\pagebreak
\section{Problem 2}
    Using \textbf{only} the definitions of $\lim$, $\limsup$, and $\liminf$, prove the following statements.
    \begin{enumerate}[label=(\alph*)]
        \item 
        $\displaystyle{\lim_{n \to \infty} \left[\sqrt{n + 1} - \sqrt{n}\right] = 0}.$ \quad \textit{Hint:} Use difference of squares.

        \item 
        $\displaystyle{\limsup_{n \to \infty} \; (-2)^n = \infty}$.

        \item 
        $\displaystyle{\limsup_{n \to \infty} \; \left[-\sqrt{n}\right]} = -\infty.$
        

        \item 
        $\displaystyle{\liminf_{n \to \infty} \left[\frac{1}{n} + (1 + (-1)^n)n\right] = 0}.$
        
    \end{enumerate}
    % \pagebreak
    \subsection{Solution (a)}
        First we prove that for all $m, n \in \mathbb{N}$, if $m \geq n$, then $\frac{1}{2\sqrt{m}} < \frac{1}{2\sqrt{n}}$.
        \begin{proof}
            To begin, we know that $n \leq m$ for $m, n \in \mathbb{N}$.
            Take the square root of both sides to find that $\sqrt{n} \leq \sqrt{m}$.
            Double both sides to find $2 \sqrt{n} \leq 2 \sqrt{m}$.
            Divide each side by what is on the other side to find $\frac{1}{2 \sqrt{m}} \leq \frac{1}{2 \sqrt{n}}$.
        \end{proof}

        By a similar logic, we prove that for all $n \in \mathbb{N}$, $\frac{1}{\sqrt{n + 1} + \sqrt{n}} < \frac{1}{2\sqrt{n}}$.
        \begin{proof}
            To begin, we know that $n < n + 1$ for $n \in \mathbb{N}$.
            Take the square root of both sides to find that $\sqrt{n} < \sqrt{n + 1}$.
            Add $\sqrt{n}$ to both sides to find $2 \sqrt{n} < \sqrt{n + 1} + \sqrt{n}$.
            Divide each side by what is on the other side to find $\frac{1}{\sqrt{n + 1} + \sqrt{n}} < \frac{1}{2\sqrt{n}}$.
        \end{proof}

        \begin{proof}
            We can start by rewriting the term within the limit.
            \begin{equation}
                \sqrt{n + 1} - \sqrt{n} = (\sqrt{n + 1} - \sqrt{n}) \frac{\sqrt{n + 1} + \sqrt{n}}{\sqrt{n + 1} + \sqrt{n}}
                    =   \frac{n + 1 - n}{\sqrt{n + 1} + \sqrt{n}}
                    =   \frac{1}{\sqrt{n + 1} + \sqrt{n}}
            \end{equation}

            Let $\varepsilon > 0$.
            Choose a natural number $N$ such that $N > \frac{1}{4\varepsilon^2}$ and let $n \in \mathbb{N}$ such that $n \geq N$.
            Take our inequality and multiply by $\varepsilon^2$ while dividing by $N$ to find that $\varepsilon^2 > \frac{1}{4N}$.
            Take the square root of both sides to find $\varepsilon > \frac{1}{2\sqrt{N}}$.
            Based on earlier statements and recalling $n \geq N$, we find $\frac{1}{2\sqrt{N}} \geq \frac{1}{2\sqrt{n}}$.
            Also based on earlier, $\frac{1}{2\sqrt{n}} > \frac{1}{\sqrt{n + 1} + \sqrt{n}}$.
            Transistively, this means $\varepsilon > \frac{1}{\sqrt{n + 1} + \sqrt{n}}$.

            Since $n \in \mathbb{N}$, $n > 0$, and the same can be said about $n + 1$.
            Therefore, $\sqrt{n} > 0$ and $\sqrt{n + 1} > 0$.
            Collectively adding, that means $\sqrt{n + 1} + \sqrt{n} > 0$.
            Multiplying everything by $(\sqrt{n + 1} + \sqrt{n})^{-1}$ twice, we find that $\frac{1}{\sqrt{n + 1} + \sqrt{n}} > 0$.
            Therefore, $\abs{\frac{1}{\sqrt{n + 1} + \sqrt{n}}} = \frac{1}{\sqrt{n + 1} + \sqrt{n}}$.

            We can therefore say that $\varepsilon > \abs{\frac{1}{\sqrt{n + 1} + \sqrt{n}}}$ for any $\varepsilon > 0$ and for any $n \geq N$ for some $N \in \mathbb{N}$.
            Therefore, $\lim_{n \to \infty} \frac{1}{\sqrt{n + 1} + \sqrt{n}} = 0$.
            Since $\frac{1}{\sqrt{n + 1} + \sqrt{n}} = \sqrt{n + 1} - \sqrt{n}$, $\lim_{n \to \infty} \left[ \sqrt{n + 1} - \sqrt{n} \right] = \lim_{n \to \infty} \frac{1}{\sqrt{n + 1} + \sqrt{n}}$.
            Ergo, $\lim_{n \to \infty} \left[ \sqrt{n + 1} - \sqrt{n} \right] = 0$.
        \end{proof}
    \pagebreak
    \subsection{Solution (b)}
        \begin{proof}
            $\limsup_{n \to \infty} (-2)^n$ is defined as $\lim_{k \to \infty} \left[ \sup_{n \geq k} (-2)^n \right]$.
            The definition of convergence to $\infty$ is that for all $M > 0$, there exists some $N \in \mathbb{N}$ so that for all $n \geq N$, $s_n > M$.

            Let $C_M = \left\lceil M \right\rceil \geq M$.
            If $C_M$ is odd, operate further on $C_M + 1$.
            
            $(-2)^n = (-1)^n \cdot 2^n$, so if $n$ is even, then $(-1)^n = 1$, so $(-2)^n = 2^n$ when $n$ is even.
            Without loss of generality (between $C_M$ when $C_M$ is even and $C_M + 1$ when $C_M$ is odd), we furthermore know that $2^{C_M} > C_M$.
            Therefore, we know that $2^{C_M} > M$.

            Therefore, $\{(-2)^n\}$ is not bounded above.
            Since $2^{n + 1} > 2^n$, and by extension $2^{n + 2} > 2^n$ and $(-2)^{n + 2} > (-2)^n$ for even $n$, the sequence is monotonically increasing.
            Therefore, $\limsup_{n \to \infty} (-2)^n = \infty$.
        \end{proof}
    % \pagebreak
    \subsection{Solution (c)}
        \begin{proof}
            By the definition of $\limsup$, $\limsup_{n \to \infty} \left[-\sqrt{n}\right] = \lim_{k \to \infty} \sup_{n \geq k} \left[-\sqrt{n}\right]$.
            By the definition of $\sup$, for all $n \geq k$, $\sup \left[-\sqrt{n}\right] \geq -\sqrt{n}$.

            Let $M < 0$.
            Let $N = \left\lceil M^2 \right\rceil$.
            Since $M < 0$, $M^2 > \abs{M}$.
            Therefore, $\left\lceil M^2 \right\rceil \geq \abs{M}$, so $N > M^2$.
            Therefore, if $n \geq N$, $n > M^2$, so $-\sqrt{n} < -\sqrt{M^2} = M$ (because $M < 0$).

            Therefore, for all $M < 0$, there exists $N \in \mathbb{N}$ such that for all $n \geq N$, $-\sqrt{n} < M$.
            Ergo, $\limsup_{n \to \infty} \left[-\sqrt{n}\right] = -\infty$.
        \end{proof}
    \pagebreak
    \subsection{Solution (d)}
        \begin{proof}
            By the definition of $\liminf$, 
            \[ \liminf_{n \to \infty} \left[\frac{1}{n} + (1 + (-1)^n)n\right] = \lim_{k \to \infty} \left[ \inf_{n \geq k} \left[\frac{1}{n} + (1 + (-1)^n)n\right] \right] \]

            Look at only the odd terms.
            If $n$ is odd, then $(1 + (-1)^n)n = (1 - 1)n = 0n = 0$, so equivalently $\frac{1}{n} + (1 + (-1)^n)n = \frac{1}{n}$.
            By the Archimedean Property, for any $\varepsilon > 0$, there are infinitely many (odd and even) $n$ such that $\varepsilon > \frac{1}{n}$.
            Also, since $\frac{1}{n} > 0$, $\frac{1}{n} > -\varepsilon$, so $0$ is the greatest lower bound of $\frac{1}{n}$.
            Therefore, $\liminf_{n \to \infty} \left[\frac{1}{n} + (1 + (-1)^n)n\right] = 0$.
        \end{proof}

\pagebreak
\section{Problem 3}
    Let $\{s_n\}$ be an arbitrary sequence in $\mathbb{R}$.
    \begin{enumerate}[label=(\alph*)]
        \item 
        Let $s \in \mathbb{R}$. Show that $\displaystyle{\limsup_{n \to \infty} s_n = s}$ if and only if the following two conditions hold:
        \begin{itemize}
            \item 
            For all $\epsilon > 0$, there is $N \in \mathbb{N}$ so that $s_n < s + \epsilon$ for all $n \geq N$.

            \item 
            For all $\epsilon > 0$ and for all $N \in \mathbb{N}$, there exists $n \geq N$ so that $s_n > s - \epsilon$.
        \end{itemize}
        
        \item 
        Show that some subsequence of $\{s_n\}$ converges (in the sense of extended real numbers, if necessary) to $\displaystyle{\limsup_{n \to \infty} s_n}$ and some subsequence of $\{s_n\}$ converges to $\displaystyle{\liminf_{n \to \infty} s_n}$.

        \item 
        Define
        \begin{equation*}
            E = \{x \in [-\infty, \infty] : \text{ some subsequence } \{s_{n_k}\} \text{ of } \{s_n\} \text{ converges to } x\}.
        \end{equation*}
        Show that $\displaystyle{\limsup_{n \to \infty} s_n = \sup E}$ and $\displaystyle{\liminf_{n \to \infty} s_n = \inf E}$.\\[5pt] \textit{Hint:} what happens to $\limsup$ and $\liminf$ on subsequences?\\
        
    \end{enumerate}
    % \pagebreak
    \subsection{Solution (a)}
        Assume that we are working within the real numbers.
        First we prove the forward direction, starting with the first statement.
        \begin{proof}
            Suppose that $\limsup_{n \to \infty} s_n = s$.
            Recall that $\limsup_{n \to \infty} s_n = \lim_{k \to \infty} \left[ \sup_{n > k} s_n \right]$.
            Now recall the definition of the limit: for all $\varepsilon > 0$, there exists some $N \in \mathbb{N}$ such that for all $n \geq N$, $d(\sup\{s_m\}_{m \geq n}, s) < \varepsilon$.
            In the standard metric, the latter term can be turned into $\abs{\sup\{s_m\}_{m \geq n} - s} < \varepsilon$.
            
            Recall the meaning of $\sup\{s_m\}_{m \geq n}$: that for all $m \geq n$, $s_m \leq \sup\{s_m\}_{m \geq n}$.
            Let $\sup\{s_m\}_{m \geq n} = \sigma$, so for all $m \geq n$, $s_m \leq \sigma$.
            Since $n \geq N$, $s_m \leq \sigma$ for all $m \geq N$.
            This means that $\sigma - s_m \geq 0$, so $\sigma - s_m = \abs{\sigma - s_m}$.

            We also know that $\abs{\sup\{s_m\}_{m \geq n} - s} < \varepsilon$, or equivalently that $\abs{\sigma - s} < \varepsilon$.
            By the absolute value, $-\varepsilon < \sigma - s < \varepsilon$.
            Add $s$ to all sides to find that $s - \varepsilon < \sigma < s + \varepsilon$.
            Earlier we established that for all $m \geq N$, $s_m \leq \sigma$, so we can transistively say for all $m \geq N$, $s_m < s + \varepsilon$.

            We earlier established $N \in \mathbb{N}$ as arbitrary (``There exists some''), and the only requirement of $\varepsilon$ is that $\varepsilon > 0$, so we can replace $\varepsilon = \epsilon$ and $m = n$.
            Ergo, for all $\epsilon > 0$, there exists $N \in \mathbb{N}$ such that for all $n \geq N$, $s_n < s + \epsilon$.
        \end{proof}

        Now we prove the second statement.
        \begin{proof}
            As with earlier, suppose that $\limsup_{n \to \infty} s_n = s$.
            For all $\varepsilon > 0$, there exists some $N \in \mathbb{N}$ such that for all $m \geq N$, $\abs{\sup\{s_k\}_{k \geq m} - s} < \varepsilon$.

            Recall a theorem we proved on the homework about $\sup\{s_k\}_{k \geq m} = \sigma$.
            We proved that for any $\epsilon > 0$, there exists some $s_n \in \{s_k\}_{k \geq m}$ where $n \geq m$ such that $s_n > \sigma - \epsilon$.
            We said in class that $\sup\{s_k\}_{k \geq m}$ is monotonically decreasing, so $\sup\{s_k\}_{k \geq m} \geq s$.
            Add $\epsilon$ to both sides of $s_n > \sigma - \epsilon$ and find $s_n + \epsilon > \sigma$, which transistively means $s_n + \epsilon > s$ and by extension $s_n > s - \epsilon$.

            Since this appplies for all $m \geq N$ (transistively $n > N$) and $\{m \in \mathbb{N} \Big| m > N\} \subset \mathbb{N}$, then by the fact that specific for all implies general exists, for all $m \in \mathbb{N}$, there exists some $n \geq m$ such that $s_n + \epsilon > s$.
            Ergo, renaming $m$ to $N$, for all $\epsilon > 0$ and for all $N \in \mathbb{N}$, there exists some $n \geq N$ such that $s_n > s - \epsilon$.
        \end{proof}

        Lastly, we prove the reverse.
        \begin{proof}
            Suppose that for all $\epsilon > 0$, there is $N \in \mathbb{N}$ so that $s_n < s + \epsilon$ for all $n \geq N$ and that for all $M \in \mathbb{N}$, there exists $n \geq M$ so that $s_n > s - \epsilon$.

            Begin with the former, that for all $\epsilon > 0$, there is $N \in \mathbb{N}$ so that $s_n < s + \epsilon$ for all $n \geq N$.
            By extension, for all $\frac{\epsilon}{2} > 0$, there is $N \in \mathbb{N}$ so that $s_n < s + \frac{\epsilon}{2}$ for all $n \geq N$.
            This means that $\sup\{s_n\}_{n \geq N} \leq s + \frac{\epsilon}{2}$ and by extension $\sup\{s_n\}_{n \geq N} < s + \epsilon$.
            Since if $A \subset B$ then $\sup A \leq \sup B$ and for all $m \geq N$, $\{s_n\}_{n \geq m} \subset \{s_n\}_{n \geq N}$, $\sup\{s_n\}_{n \geq m} \leq \sup\{s_n\}_{n \geq N}$.
            By extension, for all $m \geq N$, $\sup\{s_n\}_{n \geq m} < s + \epsilon$.
            Therefore, for all $\epsilon > 0$, there exists some $N \in \mathbb{N}$ such that for all $m \geq N$, $\sup\{s_n\}_{n \geq m} < s + \epsilon$.

            Now the latter, that for all $\epsilon > 0$, for all $M \in \mathbb{N}$, there exists $n \geq M$ so that $s_n > s - \epsilon$.
            Suppose we restrict $M$ so that $M \geq N$. 
            Then for all $\epsilon > 0$, there exists $N \in \mathbb{N}$ such that for all there exists $n \geq N$ so that $s_n > s - \epsilon$.
            Since we are working in the extended real numbers, $\sup \{s_n\}$ exists, and by extension $\sup \{s_n\}_{n \geq M}$ exists.
            Combining these, since $\sup \{s_n\}_{n \geq M} \geq s_n$ for all $n \geq M$, for all $M \in \mathbb{N}$, $\sup \{s_n\}_{n \geq M} > s - \epsilon$.
            By extension, for all $N \in \mathbb{N}$, for all $M \geq N$, $\sup \{s_n\}_{n \geq M} > s - \epsilon$.
            Renaming $M$ to $m$, for all $\epsilon > 0$ and for all $N \in \mathbb{N}$, for all $m \geq N$, $\sup \{s_n\}_{n \geq m} > s - \epsilon$.
            This guarantees that for all $\epsilon > 0$, there exists some $N \in \mathbb{N}$ such that for all $m \geq N$, $\sup \{s_n\}_{n \geq m} > s - \epsilon$.

            Combining the conclusions we came to in our previous two paragraphs, for all $\epsilon > 0$, there exists some $N \in \mathbb{N}$ such that for all $m \geq N$, $\sup\{s_n\}_{n \geq m} < s + \epsilon$ and $\sup \{s_n\}_{n \geq m} > s - \epsilon$.
            This means transistively that $s - \epsilon < \sup\{s_n\}_{n \geq m} < s + \epsilon$.
            Subtract $s$ form both sides to find $- \epsilon < \sup\{s_n\}_{n \geq m} - s < \epsilon$.
            By the definition of the absolute value, $\abs{\sup\{s_n\}_{n \geq m} - s} < \epsilon$.
            Summarizing, for all $\epsilon > 0$, there exists some $N \in \mathbb{N}$ such that for all $m \geq N$, $\abs{\sup\{s_n\}_{n \geq m} - s} < \epsilon$.
            Therefore, $\sup\{s_n\}_{n \geq m}$ converges to $s$.
            Ergo, $\limsup_{n \to \infty} s_n = s$.
        \end{proof}

        Last, we combine everything, which is a bit messy and tedious. 
        \begin{proof}
            We proved that if $\limsup_{n \to \infty} s_n = s$, then for all $\epsilon > 0$, there is $N \in \mathbb{N}$ so that $s_n < s + \epsilon$ for all $n \geq N$.
            We proved that if $\limsup_{n \to \infty} s_n = s$, then for all $\epsilon > 0$ and for all $N \in \mathbb{N}$, there exists $n \geq N$ so that $s_n > s - \epsilon$.
            We proved that if for all $\epsilon > 0$, there is $N \in \mathbb{N}$ so that $s_n < s + \epsilon$ for all $n \geq N$, and if for all $\epsilon > 0$ and for all $N \in \mathbb{N}$, there exists $n \geq N$ so that $s_n > s - \epsilon$, then $\limsup_{n \to \infty} s_n = s$.
            Ergo, $\limsup_{n \to \infty} s_n = s$ if an only if for all $\epsilon > 0$, there is $N \in \mathbb{N}$ so that $s_n < s + \epsilon$ for all $n \geq N$ and for all $\epsilon > 0$ and for all $N \in \mathbb{N}$, there exists $n \geq N$ so that $s_n > s - \epsilon$.
        \end{proof}
    \pagebreak
    \subsection{Solution (b)}
        This will be done first for $\limsup$, then for $\liminf$.
        \begin{proof}[Proof by Construction]
            Suppose that \seq{s} is bounded above and define $a_n = \sup_{m \geq n} s_m$.
            \seq{a} is decreasing and bounded below, so $a_n \to L$ where $L = \limsup_{n \to \infty} s_n$.
            Let's create a subsequence of \seq{s} choosing each $n_k$ where $s_{n_k} > a_{n_k} - \frac{1}{k}$ and $n_k > n_{k - 1}$.
            Since $a_{n_k} = \sup_{m \geq n_k} s_m$, and $\frac{1}{k} > 0$, this can be applied.
            Furthermore, $a_{n_k} \geq s_{n_k}$.
            $\frac{1}{k}$ converges to $0$, so $a_{n_k} - \frac{1}{k}$ converges to $L$ where $L = \limsup_{n \to \infty} s_n$.
            We also know that $a_{n_k}$ converges to $L$.
            Therefore, by the sequence squeeze theorem, since $a_{n_k} \geq s_{n_k} > a_{n_k} - \frac{1}{k}$, $s_{n_k}$ converges to $L$.
            Therefore, there is a subsequence of \seq{s} that converges to $\limsup_{n \to \infty} s_n$.
        \end{proof}

        Now for $\liminf$, which is basically identical with some tweaks.
        \begin{proof}[Proof by Construction]
            Suppose that \seq{s} is bounded below and define $a_n = \inf_{m \geq n} s_m$.
            \seq{a} is decreasing and bounded below, so $a_n \to L$ where $L = \liminf_{n \to \infty} s_n$.
            Let's create a subsequence of \seq{s} choosing each $n_k$ where $s_{n_k} < a_{n_k} + \frac{1}{k}$ and $n_k > n_{k - 1}$.
            Since $a_{n_k} = \inf_{m \geq n_k} s_m$, and $\frac{1}{k} > 0$, this can be applied.
            Furthermore, $a_{n_k} \leq s_{n_k}$.
            $\frac{1}{k}$ converges to $0$, so $a_{n_k} + \frac{1}{k}$ converges to $L$ where $L = \limsup_{n \to \infty} s_n$.
            We also know that $a_{n_k}$ converges to $L$.
            Therefore, by the sequence squeeze theorem, since $a_{n_k} \leq s_{n_k} < a_{n_k} - \frac{1}{k}$, $s_{n_k}$ converges to $L$.
            Therefore, there is a subsequence of \seq{s} that converges to $\liminf_{n \to \infty} s_n$.
        \end{proof}

    % \pagebreak
    \subsection{Solution (c)}
        First $\liminf$.
        \begin{proof}
            When you pass a subsequence to an infimum, you essentially pass an ordered subset to the infimum function.
            Since the ordered subset is itself a subset and the infimum of a subset will be greater than or equal to the infimum of the original set, the infimum of the subsequence will be greater than or equal to the infimum of the original sequence, so $\inf\seq{s} \leq \inf\seq[n_k]{s}$.
            Take subsets over the same span and find $\inf_{n \geq m}\seq{s} \leq \inf_{n_k \geq m}\seq[n_k]{s}$.
            Take the limit of this and we find that for every subsequence \seq[n_k]{s}, $\lim_{n \to \infty} \inf_{n \geq m}\seq{s} \leq \lim_{n \to \infty} \inf_{n_k \geq m}\seq[n_k]{s}$.
            Therefore, $\liminf_{n \to \infty}\seq{s} \leq \liminf_{n \to \infty}\seq[n_k]{s}$.
            Since it is known that if the limit exists, then the limit infimum is equal to the limit, for any \seq[n_k]{s}, $\liminf_{n \to \infty}\seq{s} \leq \lim_{n \to \infty}\seq[n_k]{s}$.
            Therefore, for any element $x \in E$, $\liminf_{n \to \infty}\seq{s} \leq x$.
            From part (b), we know that there is a subsequence of \seq{s} that converges to $\liminf_{n \to \infty}\seq{s}$.
            Therefore, $\liminf_{n \to \infty}\seq{s} \in E$.
            Earlier in this docuent we proved that if for all $x \in X$, $y \leq x$ and $y \in X$, then $y = \inf X$.
            Ergo by that proposition, $\liminf_{n \to \infty}\seq{s} = \inf E$.
        \end{proof}

\pagebreak
\section{Problem 4}
    Let $a > 1$. In this problem, we will define what it means to write $a^r$ for a positive \textit{irrational} number $r$. We first define exponentiation for positive \textit{integer} exponents in the usual way, via repeated multiplication. We know this satisfies some properties, for example
    \begin{alignat*}{2}
        a^{r + s} &= a^r a^s \qquad \qquad \qquad &&(\ast)\\[5pt]
        (ab)^{r} &= a^r b^r \qquad \qquad \qquad &&(\Delta)
    \end{alignat*}
    for $a, b > 0$ and $r, s > 0$ integers. Now, we proved in class that for any positive integer $q$, there is a \textit{unique} positive number whose $q$th power is $a$, which we denote $\sqrt[q]{a}$. For a positive rational number $r = \frac{p}{q}$, we define $a^r = \sqrt[q]{a^p}$.
    \begin{enumerate}[label=(\alph*)]
        \item 
        Show that this definition makes sense, i.e. is independent of the representation of $r$ as a fraction. That is, show that if rational numbers $\frac{p}{q} = \frac{r}{s}$ then $\sqrt[q]{a^p} = \sqrt[s]{a^r}$. \textit{Hint:} use lowest terms and the uniqueness of $q$th roots.

        \item 
        Show that properties $(\ast)$ and $(\Delta)$ above still hold for positive rational numbers. \textit{Hint:} find a common denominator and use uniqueness again.

        \item 
        For positive rational $r < s < t$, show that $a^r < a^s$ and that
        \begin{equation*}
            0 < a^s - a^r = a^r(a^{s - r} - 1).
        \end{equation*}
        Conclude that for any positive rational $r, s$ and any rational $t$ larger than both $r$ and $s$, we have
        \begin{equation*}
            |a^r - a^s| < a^t(a^{|r - s|} - 1).
        \end{equation*}

        \item 
        Now let $r$ be a positive \textit{irrational} number. We know there is a sequence of rational numbers $q_k$ that converge to $r$ (for example, the decimal expansion of $r$). We define
        \begin{equation*}
            a^r = \lim_{k \to \infty} a^{q_k}.
        \end{equation*}
        Show that the sequence $\{a^{q_k}\}$ is Cauchy, and conclude that this limit exists.\\ \textit{Hint:} Use part (c) and Rudin Theorem 3.20(d).
        
        \item
        Show that the choice of sequence $\{q_k\}$ is irrelevant; that is, show that if $\{p_k\}$ and $\{q_k\}$ are sequences of rational numbers converging to $r$, then 
        \begin{equation*}
            \lim_{k \to \infty} a^{p_k} = \lim_{k \to \infty} a^{q_k}.
        \end{equation*}

        \item 
        Using limit laws, prove that properties $(\ast)$ and $(\Delta)$ continue to hold for positive irrational exponents.
        
    \end{enumerate}

    \textit{Remark:} For $r < 0$ we then define $a^r = \frac{1}{a^{-r}}$, and for $0 < a < 1$ we define $a^r = \left(\frac{1}{a}\right)^{-r}$. It is then a matter of checking cases to verify that all the properties of exponentiation we want are indeed true.
    \pagebreak
    \subsection{Solution (a)}
        \textbf{Billie's Proposition:}
            If $\frac{p}{q} = \frac{r}{s}$ and $\frac{p}{q}$ is in lowest terms but $\frac{r}{s}$ is not, there exists some nonzero integer $t$ such that $r = pt$ and $s = qt$.
        \begin{proof}
            Let $\frac{p}{q} = \frac{r}{s}$ and $\frac{p}{q}$ is in lowest terms but $\frac{r}{s}$ is not.
            Multiply by to find that $ps = qr$.
            
            Divide by $p$ to find that $s = \frac{qr}{p}$.
            Since $s$ is a nonzero integer, $p$ divides $qr$.
            Since $\frac{p}{q}$ is in lowest terms, $p$ does not divide $q$, so $p$ must divide $r$.
            Therefore, there is some nonzero integer $t$ such that $p = rt$.

            Substitute in our value of $p$ into the result of the cross multiplication to find $rst = qr$.
            Divide both sides by $r$ to find that $q = st$.
        \end{proof}
        \begin{proof}
            Let $x = \sqrt[q]{a^p}$ and $y = \sqrt[s]{a^r}$, and suppose that $\frac{p}{q} = \frac{r}{s}$ and that $\frac{p}{q}$ is in lowest terms while $\frac{r}{s}$ is not.
            This would mean that $x = a^\frac{p}{q}$ and $y = a^\frac{r}{s}$.
            Take $x$ to the power of $q$ and $y$ to the power of $s$ to find $x^q = a^p$ and $y^s = a^r$.
            Take the former to the power of $r$ and the latter to the power of $p$ to find $x^{qr} = a^{pr}$ and $y^{ps} = a^{pr}$.
            Transistively, this means $x^{qr} = y^{sp}$.
            
            By Billie's Proposition, since $\frac{p}{q}$ is in lowest terms while $\frac{r}{s}$ is not, there exists some nonzero integer $t$ such that $r = pt$ and $s = qt$.
            Substitute these into our equation to find that $x^{qpt} = y^{qpt}$.
            Take the $qpt$th root of both sides to find that $\sqrt[qpt]{x^{qpt}} = \sqrt[qpt]{y^{qpt}}$
            By the uniqueness of $q$th roots, there is only one possible positive value of $\sqrt[qpt]{x^{qpt}}$ and one possible positive value of $\sqrt[qpt]{y^{qpt}}$, which are $x$ and $y$ respectively.
            
            Therefore, $x = y$.
            Ergo, $\sqrt[q]{a^p} = \sqrt[s]{a^r}$.
        \end{proof}
    % \pagebreak
    \subsection{Solution (b)}
        \subsubsection{Property ($\Delta$)}
            \begin{proof}
                Let $\frac{p}{q}$ be a rational number.
                Let $x = (ab)^\frac{p}{q}$ and let $y = a^\frac{p}{q} b^\frac{p}{q}$.
                Raise both sides of both to the power of $q$ to find $x^q = (ab)^p$ and $y^q = \left( a^\frac{p}{q} b^\frac{p}{q} \right)^q$.
                By property ($\Delta$) for integers, $x^p = a^p b^p$ and $y^p = a^p b^p$.
                Transistively, this means $x^p = y^p$, which we can apply the uniqueness of $q$th roots to to find that $x = y$.
                Therefore, $(ab)^\frac{p}{q} = a^\frac{p}{q} b^\frac{p}{q}$.
            \end{proof}

        \subsubsection{Property ($\ast$)}
            \begin{proof}
                Let $\frac{p}{q}$ and $\frac{r}{s}$ be rational numbers.
                Let $x = a^{\frac{p}{q} + \frac{r}{s}}$ and $y = a^\frac{p}{q} a^\frac{r}{s}$.
                Rewrite $x$ as $x = a^\frac{ps + qr}{qs}$.
                Take both to the $qs$ power to find that $x^{qs} = a^{ps + qr}$ and $y^{qs} = \left( a^\frac{p}{q} a^\frac{r}{s} \right)^{qs}$.
                Invoking property ($\Delta$) for integers, we find that $y^{qs} = a^{ps} a^{qr}$.
                Invoking property ($\Delta$) for integers, we find that $y^{qs} = a^{ps + qr}$.
                Transistively, this means $x^{qs} = y^{qs}$.
                Take the unique positive $qs$th root to find that $x = y$.
                Ergo, $a^{\frac{p}{q} + \frac{r}{s}} = a^\frac{p}{q} a^\frac{r}{s}$.
            \end{proof}
    \pagebreak
    \subsection{Solution (c)}
        First the first half.
        \begin{proof}
            Start with $r < s$ where $r$ and $s$ are natural numbers.
            Since $a > 1$, $a^r < a^s$.
            This means that $a^s - a^r > 0$. 
            Using properties proven in part (b), $a^s - a^r = a^r(a^{s - r} - a^{r - r}) = a^r(a^{s - r} - 1)$.
        \end{proof}

        Then the second half.
        \begin{proof}
            If $0 < r < s < t$, then $\abs{s - r} < t - r < t$.
            
            Suppose $r - s > 0$.
            This means $a^r - a^s = \abs{a^r - a^s} = 0$.
            As shown earlier, this means that $\abs{a^r - a^s} = a^s(a^{\abs{r - s}} - 1)$.
            Since $s < t$, $a^s < a^t$, so $\abs{a^r - a^s} = a^s(a^{\abs{r - s}} - 1) < a^t(a^{\abs{r - s}} - 1)$.
        \end{proof}
    % \pagebreak
    \subsection{Solution (d)}
        \begin{proof}
            Rudin Thorem 3.20(b) says that if $a > 0$ (which we know because $a > 1$), then $\lim_{n \to \infty} \sqrt[n]{a} = 1$.
            Therefore, $\lim_{n \to \infty} a^\frac{1}{n} = 1$.
            This means that for all $\varepsilon > 0$, there exists some $N \in \mathbb{N}$ such that for all $n > N$, $\abs{a^\frac{1}{n} - 1} < \varepsilon$.

            Since $a^r = \lim_{k \to \infty} a^{q_k}$, $r = \lim_{k \to \infty} q_k$, so $q_k$ is cauchy.
            Therefore, if $\abs{q_k - q_\ell} < \frac{1}{m}$ for some $m \in \mathbb{N}$, $a^{\abs{q_k - q_\ell}} < a^{\frac{1}{m}}$ and $0 < a^{\abs{q_k - q_\ell}} - 1 < a^{\frac{1}{m}} - 1$.
            Now suppose that we replace $\varepsilon$ with $\frac{\varepsilon}{a^t}$ since $a^t > 0$.
            Therefore, $a^{\abs{q_k - q_\ell}} - 1 < a^{\frac{1}{m}} - 1 < \frac{\varepsilon}{a^t}$.
            Multiply by $a^t$ to find $a^t(a^{\abs{q_k - q_\ell}} - 1) < \varepsilon$.
            Therefore, by part (c), $\abs{a^{q_k} - a^{q_\ell}} < \varepsilon$.
            Therefore, $\{a^{q_k}\}$ is Cauchy.

            According to Rudin 3.11(c), in $\mathbb{R}^k$, every Cauchy sequence converges.
            Therefore, since $\{a^{q_k}\} \in \mathbb{R}$ and is Cauchy, it converges, so $a^r$ exists.
        \end{proof}
    % \pagebreak
    % \subsection{Solution (e)}
    % \pagebreak
    % \subsection{Solution (f)}

\end{document}